\documentclass[a4paper, 11pt]{ltjsarticle}

% --- 日本語環境とフォント設定 ---
\usepackage{luatexja-fontspec}
\usepackage[ipaex]{luatexja-preset}

% --- 必要なパッケージ ---
\usepackage{graphicx}      % 画像を挿入するため
\usepackage{listings}      % コードを挿入するため
\usepackage{amsmath}       % 数式のため
\usepackage{geometry}      % 余白の設定
\usepackage{hyperref}      % リンク
\usepackage{booktabs}      % 綺麗な表のため
\usepackage{multirow}    % 表のセル結合

% --- ページレイアウト ---
\geometry{
  left=25mm,
  right=25mm,
  top=30mm,
  bottom=30mm
}

% --- タイトル情報 ---
\title{短期株価予測モデルの構築と非定常性の検証\\
       \large \textit{--- 特徴量エンジニアリングと時系列分割の反復的改善 ---}}
\author{（あなたの氏名）}
\date{\today}

% ======================================
% --- ドキュメント開始 ---
% ======================================
\begin{document}

\maketitle

% ======================================
\section{目的と初期計画}
% ======================================
本研究の目的は、ITセクターに属する国内上場企業の財務データに基づき、
3ヶ月（60営業日）後の株価パフォーマンス（上昇/下落）を予測する
機械学習モデル（XGBoost）を構築することである。
特に、田村氏の研究\cite{5}にならい、
ファンダメンタルズ指標とテクニカル指標を組み合わせた複合モデルの
有効性を検証する。

データソースはjQuants APIの無料プラン（直近2年分）に限定されるため、
当初の長期（1年後）予測から短期（3ヶ月後）予測へと計画を変更した。

% ======================================
\section{フェーズ1：初期実験と「68.5\%」の罠}
% ======================================
分析の初期段階では、データセットから欠損値（\texttt{NaN}）を持つ行を
すべて除去する、厳格な\texttt{dropna()}ルールを適用していた。

\subsection*{V2-Strict（財務 + PER/PBR）モデル}
このルールで\texttt{PER}や\texttt{PBR}を特徴量に加えたモデル（V2-Strict）を
学習させたところ、\textbf{正解率68.5\%}という見かけ上、非常に高い精度が観測された。
当初、我々はこれを「PER/PBR指標の有効性」によるものと誤認した。

\subsection*{欠陥の発見}
しかし、詳細な分析の結果、この高い精度はモデルの性能ではなく、
\textbf{データセットのバイアス}によるものであることが判明した。
\begin{itemize}
    \item \texttt{PER}や\texttt{PBR}は、赤字企業（Profit < 0）や
    債務超過企業（Equity \textless= 0）では計算できず、\texttt{NaN}となる。
    \item 厳格な\texttt{dropna()}ルールは、これらの企業をデータセットから
    \textbf{意図せず全てフィルタリング}していた。
    \item 結果として、モデルは予測が容易な**「優良企業（黒字企業）」のみを学習**
    していたため、正解率が人為的につり上がっていた。
\end{itemize}

% ======================================
\section{フェーズ2：実験計画の再立案と公正なベースラインの確立}
% ======================================
前述の欠陥を受け、本研究のゴールを
「赤字企業なども含む、より現実的な**『ごちゃまぜデータ』**全体で
正解率を上げること」に再設定した。
そのために、以下の2つの根本的な修正を加えた。

\subsection*{1. 公正なdropnaルール（賢いフィルター）}
厳格な\texttt{dropna()}を廃止し、\texttt{dropna(subset=['target', 'ROE'])}
というルールを採用した。これは、
\begin{itemize}
    \item \texttt{target}が\texttt{NaN}の「未来の行」
    \item \texttt{ROE}が\texttt{NaN}の「計算不能な（債務超過や財務データ未結合の）行」
\end{itemize}
のみを削除する、赤字企業（`Profit` < 0）は除外しない、
最も公正なフィルタリングである。

\subsection*{2. 公正な時系列分割}
当初の`len()`（行数）に基づく分割が、意図しない「銘柄間予測」
（例：楽天で学習しSHIFTでテスト）になっていた欠陥を修正。
全データを日付でソートし、時間的な**「分割日（\texttt{split\_date}）」**
を基準に学習データとテストデータを分離する、正しい時系列分割を実装した。

\subsection*{3. 公正な特徴量の定義}
「ごちゃまzeデータ」において、`Profit`や`Close`といった「生の数字」
（絶対値）は、企業規模が異なるため比較不可能（リンゴとスイカの問題）
であると判断。モデルに与える特徴量を、企業間で比較可能な
**「比率」**（共通のモノサシ）に限定した。

% ======================================
\section{フェーズ3：公正な比較実験（V1 vs V2 vs V2\_Close）}
% ======================================
上記の新ルールに基づき、ファンダメンタルズ分析の有効性を
公正に再実験した。

\subsection*{V1-FAIR（財務比率のみ）}
\begin{itemize}
    \item \textbf{特徴量:} \texttt{ROE}, \texttt{SelfCapitalRatio}
    \item \textbf{正解率: 約 50\%}
    \item \textbf{考察:} 「ごちゃまzeデータ」の短期予測において、純粋な財務比率は
    予測の役に立たない（ランダムと同等）ことが確定した。
    これが我々の**真のベースライン**である。
\end{itemize}

\subsection*{V2-FAIR（財務 + バリュー比率）}
\begin{itemize}
    \item \textbf{特徴量:} V1 + \texttt{PER}, \texttt{PBR}
    \item \textbf{正解率: 50.9\%}
    \item \textbf{考察:} V1と同様、伝統的なバリュー指標も短期予測には
    無力であることが示された。
\end{itemize}

\subsection*{V2\_Close-FAIR（財務 + 部品）}
\begin{itemize}
    \item \textbf{特徴量:} V1 + \texttt{Close}, \texttt{IssuedShares}（「生の数字」）
    \item \textbf{正解率: 56.6\%}
    \item \textbf{考察:} 「生の数字」を与えたことで、モデルが銘柄間の"さじ加減"を
    必死に学習しようと試み、正解率はわずかに改善した。
    しかし、このアプローチも50\%台の壁を越えられず、
    ファンダメンタルズ分析の限界が示された。
\end{itemize}

\subsection*{フェーズ3の結論}
ファンダメンタルズおよびバリュー指標という「第1の武器」だけでは、
「ごちゃまzeデータ」の短期予測問題は解けないことが証明された。

% ======================================
\section{フェーズ4：複合モデル（V6）と非定常性の発見}
% ======================================
ファンダメンタルズの限界を受け、田村氏の研究にならい、
性質の異なる**「第2の武器（テクニカル指標）」**を複合させた
V6モデルを構築した。

\subsection*{V6モデルの構築}
V2（財務+バリュー）の特徴量に加え、
4種類のテクニカル指標（\texttt{MAdivergence}, \texttt{RSI},
\texttt{Volatility}, \texttt{Momentum}）を、それぞれ3期間（25, 60, 75日）で
計算した、全17個の「共通のモノサシ」を特徴量として投入した。

\subsection*{ウォークフォワード分析（60-20-20分割）}
モデルの安定性を検証するため、時系列で60\%を学習、
続く20\%を「テスト1」、最後の20\%を「テスト2」として
2段階で評価した。

\subsection*{V6モデルの結果}
\begin{table}[h]
  \centering
  \caption{V6 複合モデル ウォークフォワード分析結果}
  \label{tab:v6_results}
  \begin{tabular}{lcc}
    \toprule
    評価期間 & 分割日 & 正解率 (Accuracy) \\
    \midrule
    \textbf{テスト1（平常時）} & 2024-12-24 〜 2025-03-07 & \textbf{0.5978 (59.8\%)} \\
    \textbf{テスト2（異常時）} & 2025-03-07 〜 & \textbf{0.4740 (47.4\%)} \\
    \bottomrule
  \end{tabular}
\end{table}

\subsection*{V6モデルの考察}
\begin{enumerate}
    \item \textbf{複合モデルの有効性（テスト1）:}
    「全部乗せ」の複合モデルは、**テスト1（平常時）において
    59.8\%の正解率を達成**し、V1〜V2\_Closeのベースライン（50\%〜56%）
    の壁を明確に突破した。これは、田村氏の論文の結論通り、
    **性質の異なる特徴量をアンサンブル（複合）することが短期予測に有効**
    であることを証明している。

    \item \textbf{非定常性の発見（テスト2）:}
    テスト2の期間は、市場が歴史的な急騰・急落を経験した
    「異常時」であった。モデルは「平常時」のパターンで学習していたため、
    この未知の変動に対応できず、**正解率は47.4\%へと破綻した**。

    \item \textbf{特徴量の重要度:}
    モデルは、\texttt{SelfCapitalRatio}（自己資本比率）や\texttt{ROE}といった
    **財務指標を予測の主軸**として最も重視した。テクニカル指標は、
    それらを補助する「アシスト役」として機能し、チーム全体の精度を
    引き上げたと推察される。
\end{enumerate}

\subsection*{最終結論}
本研究は、田村氏の論文（第3章・第6章）の核心的な主張を
自ら追体験・証明するものとなった。
\begin{itemize}
    \item ファンダメンタルズ分析のみでは短期予測は困難である。
    \item ファンダメンタルズとテクニカルを複合させたモデルは、
    「平常時」において予測精度を改善できる。
    \item しかし、そのモデルも市場の「非定常性」（異常時）には
    対応できないという「時系列普遍性の課題」を抱えている。
\end{itemize}

\begin{figure}[h]
  \centering
  \fbox{（ここに feature\_importance\_v6\_multi\_window.png の画像を挿入）}
  \caption{V6 複合モデルの特徴量重要度}
\end{figure}

\end{document}